\paragrafos{Θετικοί και αρνητικοί ακέραιοι: ο κόσμος κάτω από το μηδέν}{Α.2.1, «Θετικοί, αρνητικοί, ομόσημοι, ετερόσημοι ακέραιοι»}

\begin{anagke}
Η Α3 κάνει εκδρομή στο Μέτσοβο. Στο πούλμαν το ψηφιακό θερμόμετρο γράφει $-4$. «Τι είναι αυτό το μείον;» ρωτάει η Πολυξένη. «Τέσσερις βαθμοί κάτω από το μηδέν», απαντάει ο Νίκος. Ο Στέλιος ταράζεται: «Κάτω από το μηδέν; Το μηδέν δεν είναι το τίποτα; Πώς μπορεί να υπάρχει κάτι μικρότερο από το τίποτα;».

Στο ξενοδοχείο ο ανελκυστήρας έχει κουμπιά $-2$, $-1$, $0$, $1$, $2$, $3$. Ο Ηλίας πατάει το $-1$ και κατεβαίνει στο υπόγειο με το πάρκινγκ. Η Αγγελική βλέπει στον πίνακα του καιρού: Αθήνα $+7$, Ιωάννινα $-3$, Φλώρινα $-10$, Χανιά $+11$. Η Ρήνη θυμάται ότι στο ταμείο της παρέας κάποτε έλειπαν $8$ ευρώ, ένα χρέος. Και στο ντοκιμαντέρ που είδαν, ο καρχαρίας βρισκόταν $200$ μέτρα κάτω από την επιφάνεια ενώ το αεροπλάνο πετούσε $200$ μέτρα πάνω από αυτήν.

Η κυρία Αλεξάνδρα το συνοψίζει: «Οι φυσικοί αριθμοί που ξέρετε μετράνε πόσα. Όταν όμως χρειάζεται να πούμε «πάνω ή κάτω», «κέρδος ή ζημιά», «δεξιά ή αριστερά», οι φυσικοί δεν φτάνουν. Χρειαζόμαστε αριθμούς με κατεύθυνση». Κάθε «δεξιά» προϋποθέτει ένα «αριστερά» και κάθε «πάνω» ένα «κάτω»: ο κόσμος των αριθμών έχει δύο πλευρές. Σε αυτή την παράγραφο γνωρίζουμε τους θετικούς και τους αρνητικούς ακέραιους, τους βάζουμε πάνω στον άξονα και μαθαίνουμε να λέμε ποιοι έχουν το ίδιο πρόσημο και ποιοι όχι.
\end{anagke}

\begin{lexi}{θετικός, αρνητικός, ακέραιος}{θετικός, αρνητικός, ακέραιος}
\lexiword{Θετικός} είναι αυτός που «τίθεται», που βάζουμε ως βάση. \lexiword{Αρνητικός} είναι αυτός που «αρνείται», που πηγαίνει ανάποδα από τον θετικό. Το \lexiword{πρόσημο} είναι το «προ-σήμα», το σημάδι που μπαίνει πριν από τον αριθμό. Ο \lexiword{ακέραιος} είναι ο «ακέραιος», ο ολόκληρος: χωρίς κομμάτια. \lexiword{Ομόσημοι} είναι οι αριθμοί με το «ίδιο σήμα» και \lexiword{ετερόσημοι} αυτοί με «άλλο» (έτερο) σήμα. Η \lexiword{τετμημένη} είναι αυτή που «τέμνεται», που κόβουμε από την αρχή του άξονα ως το σημείο.
\end{lexi}

\begin{metafrasi}
Τέσσερις βαθμοί κάτω από το μηδέν. & $-4$\\
Αεροπλάνο $200$ μ. πάνω από τη θάλασσα. & $+200$\\
Καρχαρίας $200$ μ. κάτω από την επιφάνεια. & $-200$\\
Χρέος $8$ ευρώ. & $-8$\\
Κέρδος $15$ ευρώ. & $+15$\\
Δεύτερο υπόγειο. & $-2$\\
Δύο αριθμοί με το ίδιο πρόσημο. & ομόσημοι\\
Ο αριθμός που αντιστοιχεί σε σημείο του άξονα. & τετμημένη, $\mathrm{Α}(-3)$\\
\end{metafrasi}

\tmima{Η θεωρία}

\tmimasmall{Πού χρησιμοποιούμε αρνητικούς αριθμούς}

Οι αρνητικοί αριθμοί εκφράζουν το αντίθετο από κάτι που ήδη μετράμε με φυσικούς:
\begin{itemize}[leftmargin=1.2em,itemsep=0.3ex,topsep=2pt]
\item \textbf{Κέρδος και ζημιά.} Έσοδο, κέρδος, είσπραξη: θετικό ποσό. Έξοδο, ζημιά, χρέος: αρνητικό ποσό.
\item \textbf{Κατεύθυνση.} Δεξιά (θετική) και αριστερά (αρνητική). Πάνω (θετική) και κάτω (αρνητική). Αύξηση και μείωση.
\item \textbf{Στροφή.} Αντίθετα από τους δείκτες του ρολογιού (θετική στροφή) και κατά τη φορά των δεικτών (αρνητική στροφή).
\end{itemize}
Την ένδειξη $-2$ στο θερμόμετρο τη διαβάζουμε «πλην δύο» ή «δύο υπό το μηδέν».

\begin{orismos}{πρόσημο, θετικός, αρνητικός}
Τα σύμβολα «$+$» και «$-$», που γράφονται πριν από τους αριθμούς, λέγονται \textbf{πρόσημα}. Οι αριθμοί με πρόσημο $+$ (ή χωρίς πρόσημο) λέγονται \textbf{θετικοί}: $+5$, $+2$, $7$. Οι αριθμοί με πρόσημο $-$ λέγονται \textbf{αρνητικοί}: $-1$, $-2$, $-6$.
\end{orismos}

Το $+$ μπροστά από έναν θετικό συχνά παραλείπεται: γράφουμε $2$ αντί για $+2$. Το $-$ μπροστά από έναν αρνητικό \textbf{δεν παραλείπεται ποτέ}. Το $0$ \textbf{δεν είναι ούτε θετικός ούτε αρνητικός}.

\begin{orismos}{ομόσημοι, ετερόσημοι}
Δύο αριθμοί είναι \textbf{ομόσημοι} όταν έχουν το ίδιο πρόσημο και \textbf{ετερόσημοι} όταν έχουν διαφορετικό πρόσημο.
\end{orismos}

Για παράδειγμα οι $+2$ και $+5$ είναι ομόσημοι, οι $-1$ και $-3$ επίσης. Οι $+2$ και $-5$ είναι ετερόσημοι.

\begin{orismos}{ακέραιοι αριθμοί}
\textbf{Ακέραιοι} είναι οι φυσικοί αριθμοί μαζί με τους αντίστοιχους αρνητικούς: $$\dots,\,-4,\,-3,\,-2,\,-1,\,0,\,+1,\,+2,\,+3,\,+4,\dots$$
\end{orismos}

Οι φυσικοί $0,\,1,\,2,\,3,\dots$ είναι λοιπόν και ακέραιοι. Αν στους φυσικούς προσθέσουμε και τους αρνητικούς $-1,\,-2,\,-3,\dots$ παίρνουμε όλους τους ακέραιους. Ο $-3$ είναι ακέραιος αλλά \emph{όχι} φυσικός.

\tmimasmall{Μια ιστορική σημείωση}

Οι αρνητικοί αριθμοί εμφανίζονται για πρώτη φορά σε ένα κινέζικο μαθηματικό βιβλίο, τα «Μαθηματικά σε εννέα βιβλία», που τοποθετείται στην εποχή της δυναστείας των Χαν (206 π.Χ. -- 220 μ.Χ.). Στην ευρωπαϊκή παράδοση ο Διόφαντος (γύρω στο 250 μ.Χ.) τους χρησιμοποιεί στους ενδιάμεσους υπολογισμούς του, αλλά ως λύση ενός προβλήματος δέχεται πάντα θετικό αριθμό. Για πολλούς αιώνες οι μαθηματικοί τους θεωρούσαν «παράλογους» ή «ψεύτικους» αριθμούς. Σήμερα τους χρησιμοποιούμε στο θερμόμετρο, στον λογαριασμό της τράπεζας, στο υψόμετρο.

\tmimasmall{Οι ακέραιοι στην αριθμογραμμή: η τετμημένη}

Στον ημιάξονα $Ox$ τοποθετούσαμε τους φυσικούς δεξιά από την αρχή $\mathrm{Ο}$. Αν τον προεκτείνουμε προς τα αριστερά (στον αντικείμενο ημιάξονα $Ox'$) με την ίδια μονάδα, παίρνουμε τον \textbf{άξονα $x'x$}. Οι αρνητικοί παριστάνονται με σημεία συμμετρικά, ως προς το $\mathrm{Ο}$, των σημείων που παριστάνουν τους αντίστοιχους θετικούς. Έτσι οι θετικοί βρίσκονται δεξιά του $\mathrm{Ο}$ και οι αρνητικοί αριστερά του.

\begin{center}
\akgr[0.6]{-10}{10}{}
\end{center}

Η θέση ενός σημείου πάνω στον άξονα $x'x$ ορίζεται από έναν ακέραιο που λέγεται \textbf{τετμημένη} του σημείου. Γράφουμε $\mathrm{Α}(4)$, $\mathrm{Β}(-2)$, $\mathrm{Ο}(0)$.

\begin{center}
\akgr[1]{-5}{5}{-2/\mathrm{Β},0/\mathrm{Ο},4/\mathrm{Α}}
\end{center}

Το σημείο $\mathrm{Α}$ έχει τετμημένη $4$, το $\mathrm{Β}$ έχει $-2$ και το $\mathrm{Ο}$ έχει $0$. Όσο πηγαίνουμε προς τα αριστερά οι αριθμοί μικραίνουν.

Ο ίδιος άξονας μπορεί να σταθεί και κατακόρυφα, όπως στον ανελκυστήρα ή στο θερμόμετρο:

\begin{center}
\begin{tikzpicture}[y=0.55cm]
  \draw[-{Stealth},thick] (0,-2.7)--(0,3.9);
  \foreach \j in {-2,-1,0,1,2,3}{\draw (-0.1,\j)--(0.1,\j) node[right=3pt,font=\small]{$\j$};}
  \node[left=4pt,font=\small] at (0,3) {τρίτος όροφος};
  \node[left=4pt,font=\small] at (0,0) {ισόγειο};
  \node[left=4pt,font=\small] at (0,-2) {δεύτερο υπόγειο};
\end{tikzpicture}
\end{center}

\begin{prosoxi}
\begin{itemize}[leftmargin=1.2em,itemsep=0.3ex,topsep=2pt]
\item Το $-3$ δεν είναι «τρία ευρώ λιγότερο από το τίποτα». Είναι μια θέση, τρεις μονάδες αριστερά του $0$, και η μονάδα μπορεί να είναι βαθμός, ευρώ ή μέτρο.
\item Το $0$ δεν έχει πρόσημο. Δεν γράφουμε $+0$ ή $-0$.
\item Όλοι οι αρνητικοί είναι μικρότεροι από το $0$. Το $-100$ δεν είναι «μεγάλος αριθμός»: είναι πολύ αριστερά.
\item Μην μπερδεύεις το πρόσημο με την πράξη: στο $5-3$ το μείον είναι αφαίρεση, στο $-3$ είναι πρόσημο.
\item Οι αριθμοί $+5$ και $5$ είναι ο ίδιος αριθμός. Αντίθετα, οι $5$ και $-5$ είναι διαφορετικοί.
\end{itemize}
\end{prosoxi}

\tmima{Λυμένα παραδείγματα}

\paradeigma{Αριθμοί για μεγέθη}
\tixero Αθήνα $7$° πάνω από το $0$, Ιωάννινα $3$° κάτω, Φλώρινα $10$° κάτω. Ένα χρέος $12$ ευρώ. Μια αύξηση $200$ ευρώ και μια μείωση $50$ ευρώ.\par
\tithelo Να τα γράψω ως ακέραιους.\par
\possindeo Πάνω από το μηδέν, αύξηση, κέρδος είναι θετικό. Κάτω από το μηδέν, μείωση, χρέος είναι αρνητικό.\par
\apantisi Αθήνα $+7$, Ιωάννινα $-3$, Φλώρινα $-10$, χρέος $-12$, αύξηση $+200$, μείωση $-50$.\par
\elegxos Η Φλώρινα ($-10$) είναι πιο κρύα από τα Ιωάννινα ($-3$): θα το συγκρίνουμε στην επόμενη παράγραφο.

\paradeigma{Θετικοί, αρνητικοί, ομόσημοι, ετερόσημοι}
\tixero Οι αριθμοί $+4$, $-9$, $-1$, $+7$.\par
\tithelo Ποιοι είναι θετικοί, ποιοι αρνητικοί και ποια ζεύγη είναι ομόσημα ή ετερόσημα.\par
\possindeo Κοιτάζω μόνο τα πρόσημα. Θετικοί: $+4$ και $+7$. Αρνητικοί: $-9$ και $-1$. Ομόσημα ζεύγη είναι αυτά που έχουν το ίδιο πρόσημο, όπως $(+4,\,+7)$ και $(-9,\,-1)$. Τα υπόλοιπα, π.χ. $(+4,\,-9)$, είναι ετερόσημα.\par
\apantisi Ομόσημοι: $+4,\,+7$ και $-9,\,-1$. Ετερόσημοι: κάθε θετικός με κάθε αρνητικό.\par
\elegxos Δύο αριθμοί με το ίδιο πρόσημο βρίσκονται στην ίδια πλευρά του $0$ στον άξονα.

\paradeigma{Φυσικοί ή ακέραιοι;}
\tixero Οι αριθμοί $-4$, $0$, $15$, $-1$, $3$, $-100$.\par
\tithelo Ποιοι είναι φυσικοί και ποιοι ακέραιοι.\par
\possindeo Φυσικοί είναι οι $0,\,1,\,2,\dots$ Οι αρνητικοί δεν είναι φυσικοί. Όλοι οι αριθμοί της λίστας όμως είναι ακέραιοι.\par
\apantisi Φυσικοί: $0$, $15$, $3$. Ακέραιοι: και οι έξι.\par
\elegxos Κάθε φυσικός είναι ακέραιος. Το αντίστροφο δεν ισχύει: ο $-4$ είναι ακέραιος αλλά όχι φυσικός.

\paradeigma{Σημεία στον άξονα}
\tixero Ο άξονας του σχήματος και τα σημεία $\mathrm{Δ}$, $\mathrm{Ε}$, $\mathrm{Ζ}$.\par
\tithelo Να βρω τις τετμημένες τους και να τοποθετήσω τον $-6$.\par
\possindeo Μετράω μονάδες από το $0$, δεξιά με $+$ και αριστερά με $-$.\par
\begin{center}
\akgr[0.7]{-8}{8}{-5/\mathrm{Δ},1/\mathrm{Ε},6/\mathrm{Ζ}}
\end{center}
\apantisi $\mathrm{Δ}(-5)$, $\mathrm{Ε}(1)$, $\mathrm{Ζ}(6)$. Ο $-6$ βρίσκεται $6$ μονάδες αριστερά του $0$, δηλαδή μία μονάδα πιο αριστερά από το $\mathrm{Δ}$.\par
\elegxos Τα $\mathrm{Δ}$ και $\mathrm{Ε}$ είναι σε διαφορετικές πλευρές του $0$: η τετμημένη του $\mathrm{Δ}$ είναι αρνητική και του $\mathrm{Ε}$ θετική.

\tmima{Ασκήσεις}

\zoni{Το κατάλαβα}{πρόσημα, ακέραιοι, άξονας}

\ex Εξέφρασε με ακέραιο: α) $5$ βαθμοί κάτω από το μηδέν \quad β) $12$ μ. πάνω από τη θάλασσα \quad γ) χρέος $8$ € \quad δ) κέρδος $15$ € \quad ε) το δεύτερο υπόγειο \quad στ) $30$ μ. κάτω από την επιφάνεια της θάλασσας.

\ex Ποιοι από τους αριθμούς $-7$, $+3$, $0$, $-12$, $8$, $+1$ είναι θετικοί, ποιοι αρνητικοί και ποιος δεν είναι ούτε το ένα ούτε το άλλο;

\ex Πες αν είναι ομόσημοι ή ετερόσημοι: α) $-3$ και $-8$ \quad β) $+4$ και $-4$ \quad γ) $5$ και $+9$ \quad δ) $-2$ και $6$ \quad ε) $-1$ και $-100$ \quad στ) $+8$ και $-8$.

\ex Η Ρήνη λέει «Αντί για $+2$ μπορώ να γράψω $2$», και ο Στέλιος «Αντί για $-2$ μπορώ να γράψω $2$». Ποιος έχει δίκιο;

\ex Σχεδίασε άξονα με μονάδα $1$ cm και σημείωσε τους $-4$, $-1$, $0$, $2$, $5$.

\ex Στον άξονα βρες τις τετμημένες των $\mathrm{Α},\,\mathrm{Β},\,\Gamma,\,\mathrm{Δ}$.
\begin{center}
\akgr{-6}{6}{-4/\mathrm{Α},-1/\mathrm{Β},3/\Gamma,5/\mathrm{Δ}}
\end{center}

\ex Σωστό ή λάθος; α) Το $0$ είναι θετικός αριθμός. \quad β) Κάθε φυσικός αριθμός είναι ακέραιος. \quad γ) Κάθε ακέραιος είναι φυσικός. \quad δ) Οι $-3$ και $-5$ είναι ετερόσημοι. \quad ε) Ο $-2$ είναι φυσικός αριθμός. \quad στ) Στον άξονα οι αρνητικοί βρίσκονται αριστερά του $0$.

\zoni{Το χρησιμοποιώ}{θερμόμετρα, ασανσέρ, θάλασσες}

\ex \textbf{Ο καιρός.} Οι προβλέψεις: Αθήνα $7$° πάνω από το $0$, Θεσσαλονίκη $3$° κάτω, Ιωάννινα $5$° κάτω, Χανιά $11$° πάνω, Φλώρινα $10$° κάτω. Γράψε τις θερμοκρασίες ως ακέραιους και τοποθέτησέ τες σε έναν κατακόρυφο άξονα.

\ex \textbf{Ο ανελκυστήρας.} Οι όροφοι του ξενοδοχείου είναι $-2,\,-1,\,0,\,1,\,2,\,3$. Ο Βασίλης είναι στον $3$ και θέλει το δεύτερο υπόγειο. Πόσους ορόφους κατεβαίνει;

\ex \textbf{Βουνό και θάλασσα.} Το αεροπλάνο πετάει στα $+300$ μ., το ψάρι στα $-45$ μ. και το υποβρύχιο στα $-120$ μ. Ποιο από τα τρία είναι πιο κοντά στην επιφάνεια της θάλασσας;

\ex \textbf{Το ταμείο.} Η Ρήνη έβαλε $20$ €, ο Στέλιος πήρε $7$ €, η Φανή έβαλε $5$ € και ο Βασίλης πήρε $12$ €. Γράψε κάθε κίνηση ως ακέραιο.

\ex \textbf{Ο καιρός αλλάζει.} Το πρωί στο Μέτσοβο είναι $-3$° και μέχρι το μεσημέρι ανεβαίνει $5$ βαθμούς. Βρες τη νέα θερμοκρασία μετρώντας πάνω στον άξονα.

\zoni{Το σκέφτομαι}{είναι το μηδέν τίποτα;}

\ex Ο Στέλιος ισχυρίζεται ότι δεν υπάρχει αριθμός μικρότερος από το $0$. Φέρε αντιπαράδειγμα από τη ζωή.

\ex Γιατί το $0$ δεν είναι ούτε θετικός ούτε αρνητικός αριθμός;

\ex Τα $0$°C είναι η θερμοκρασία πήξης του νερού και τα $32$°F είναι η ίδια θερμοκρασία σε βαθμούς Φαρενάιτ. Τι δείχνει αυτό για το «μηδέν» στο θερμόμετρο;

\ex Πόσοι ακέραιοι υπάρχουν από το $-5$ ως το $5$, με τα άκρα; Πόσοι ανάμεσα στο $-5$ και στο $5$ χωρίς τα άκρα;

\zoni{Για το σπίτι}{ψάξε τους αρνητικούς γύρω σου}

\ex Βρες πέντε περιπτώσεις της καθημερινής ζωής που χρησιμοποιούνται αρνητικοί αριθμοί.

\ex Μέτρησε τη θερμοκρασία κάθε πρωί για μια εβδομάδα και παράστησέ τη σε άξονα.
